\scanpage{85}
- \mspage{IV-1013} -
Comme par hypothèse la suite $\underline f$ est $\operatorname{gr}_{\underline K}^\bullet(A)$-régulière, elle est aussi
$\operatorname{gr}_{\underline K}^\bullet(N)$-régulière; on peut par suite appliquer (19.6.3), car tout module
quotient d'un sous-module d'un \deleted{\illegible{1}}
$\operatorname{gr}_{\underline K}^k(M)$ est alors séparé pour
la topologie $\underline J$-préadique puisque les $f_i$ appartiennent au radical de $A$
par hypothèse. Le diagramme commutatif (19.6.2.1) s'écrit donc ici

(19.7.1.1)
\diagram{\begin{tikzcd}[ampersand replacement=\&,column sep=large]
\&0\arrow[d]\\
(\operatorname{gr}_{\underline K}^\bullet(R,N)\otimes_{A/\underline K}(A/\underline L))[T_1,\ldots,T_n]
 \arrow[r,"\psi'_R"]\arrow[d]
\&\operatorname{gr}_{\underline L}^\bullet(R,N)\arrow[d]\\
(\operatorname{gr}_{\underline K}^\bullet(N)\otimes_{A/\underline K}(A/\underline L))[T_1,\ldots,T_n]
 \arrow[r,"\psi_N"]\arrow[d]
\&\operatorname{gr}_{\underline L}^\bullet(N)\arrow[d]\\
(\operatorname{gr}_{\underline K}^\bullet(\unsure{M|N})\otimes_{A/\underline K}(A/\underline L))[T_1,\ldots,T_n]
 \arrow[r,"\psi_M"]\arrow[d]
\&\operatorname{gr}_{\underline L}^\bullet(M)\arrow[d]\\
0\&0
\end{tikzcd}}
où $\operatorname{gr}_{\underline K}^\bullet(R,N)$ \deleted{\illegible{2}}
(resp. $\operatorname{gr}_{\underline L}^\bullet(R,N)$) est le gradué associé à $R$
pour la filtration des $R\cap\underline K^mN$ (resp. \deleted{\unsure{des}} $R\cap\underline L^mN$);
rappelons que dans
ce diagramme les colonnes sont exactes et que $\psi_N$ est \uline{bijectif}. Il s'agit,
en vertu de (19.6.3), de prouver que $\psi'_R$ est \uline{surjectif}, ce qui
va être \deleted{\unsure{prouvé}}\inserted{fait} en plusieurs étapes.

(19.7.1.2) Posons $B=\operatorname{gr}_{\underline K}^\bullet A\otimes_{A/\underline K}(A/\underline L)$,
$P=(\operatorname{gr}_{\underline K}^\bullet(N)\otimes_{A/\underline K}(A/\underline L)$,
$Q=\operatorname{gr}_{\underline L}^\bullet(R,N)$;
puisque $\psi_N$ est bijectif, \deleted{\illegible{1}} $P[T_1,\ldots,T_n]$ s'identifie à
$\operatorname{gr}_{\underline L}^\bullet(N)$, et
\deleted{\unsure{identifie}} $Q$ à un sous-$B[T_1,\ldots,T_n]$-module de $P[T_1,\ldots,T_n]$;
nous allons
d'abord voir que l'on a
\[
\text{(19.7.1.3)}\qquad Q=(Q\cap P)[T_1,\ldots,T_n].
\]
Pour cela, posons $Z=Q/((Q\cap P)[T_1,\ldots,T_n])$; c'est un sous-$B[T_1,\ldots,T_n]$-
\mspage{84}
